% KOMA-Script's footnotes in Times: scrartcl's default \deffootnote, whose
% text starts left of its mark. Each footnote keeps the words and spaces of
% the PDF, whatever its first word: a narrow one ("I", "a"), a wide one,
% one in italics, sans-serif, or typewriter, and one with punctuation.
% Also a footnote with its own number, one in a minipage, and \thanks.
\documentclass[fontsize=11]{scrartcl}
\usepackage{times}
\usepackage[T1]{fontenc}
\title{Spaced notes\thanks{I thank a funder of this work.}}
\author{Someone}
\begin{document}
\maketitle
Text one.\footnote{I thank a referee for this.}
Text two.\footnote{A measure is finite here.}
Text three.\footnote{We thank the editor too.}
Text four.\footnote{I am grateful to her.}
Text five.\footnote[20]{I owe this point to him.}
Text six.\footnote{\emph{I} stress this word.}
Text seven.\footnote{I, too, agree with it.}
Text eight.\footnote{Ian thanks a colleague.}
Text nine.\footnote{a fortiori it holds.}
Text ten.\footnote{\textsf{I} say this in sans.}
Text eleven.\footnote{\texttt{I} type this in mono.}
Text twelve.\footnote{I.e. the same again.}

\begin{minipage}{0.8\textwidth}
In a box.\footnote{I note this in a box.}
\end{minipage}
\end{document}
